\section{Related Work}
\label{sec:related}

\paragraph{Simultaneous speech translation.} Wait-$k$~\citep{ma2019stacl} writes after a fixed delay;
LocalAgreement~\citep{liu2020low,polak2022cuni,machacek2023turning} commits the prefix two consecutive hypotheses agree
on; learned policies emit a wait token~\citep{koshkin2024transllama} or read the decoder's attention to the
audio~\citep{papi2023edatt,papi2024streamatt}. Latency is measured with SimulEval~\citep{ma2020simuleval} and LAAL~\citep{papi2022over};
pre-segmented evaluation overstates deployability~\citep{papi2025real}. SimulS2ST systems
keep one speaker's voice~\citep{zhang2024streamspeech,seamless2023,labiausse2025hibiki,cheng2025liveinterpret,he2026simuls2st,deng2025simuls2s}.

\paragraph{Latent and learned timing.} Monotonic attention marginalises over write positions
(MILk;~\citealp{arivazhagan2019milk}), also per attention head (MMA;~\citealp{ma2020mma}) and for streaming speech (EMMA,
the policy of SeamlessStreaming;~\citealp{ma2023emma}); all add a latency loss, because likelihood alone prefers late
writes. The drift of our write posterior (\S\ref{sec:write}) is this known effect. The Hidden Markov Transformer
treats the moment a target token starts as a hidden state~\citep{zhang2023hmt}, and Hibiki places each target word after the source that makes it
predictable~\citep{labiausse2025hibiki}. Our write alignment differs in its state space (written chunks of several
speakers and their order) and, like MILk, needs a latency term. These policies time the \emph{text}; for speech output we time
\emph{playback}, where the state is the queue of synthesised speech, not the source read so far. Posterior
regularisation constrains expectations under a latent posterior by a KL projection~\citep{ganchev2010posterior}; we use
it to bound the delay at which synthesised words are heard. Automatic dubbing controls the length of the
translation~\citep{lakew2022isometric}; our playback layer instead chooses among renderings of a fixed text, by their
effect on the speaker's backlog.

\paragraph{Multi-speaker recognition and translation.} Serialised output training transcribes overlapping speakers
in order~\citep{kanda2020sot,kanda2021saasr}, token-level serialisation does so in streaming~\citep{kanda2022tsot} and
also for joint streaming recognition and translation~\citep{papi2023tsot}; streaming diarizers keep speaker order over
long recordings~\citep{park2024sortformer,medennikov2025streaming}, LLMs can post-edit diarized
transcripts~\citep{wang2024diarizationlm}, and diarization can condition a recogniser~\citep{polok2026dicow,huo2026tagspeech}
or be interleaved with a speech LLM's audio tokens (G-STAR;~\citealp{peng2026gstar}), which we adopt. STAC-ST translates
two-party calls offline with turn tokens~\citep{zuluaga2023stac}; DiariST~\citep{yang2024diarist} streams
speaker-attributed translation with token-level serialisation; X-Translator~\citep{zhao2026xtranslator} binds each
committed recognition span to the speaker with the largest temporal overlap. Voice-cloned translation lowers listeners'
workload in online meetings~\citep{ma2026samevoice}.
